% =============================================================================
% robo/emi.tex — Documento de prueba
% Contenido: Tesis EMI "Plataforma de Mantenimiento Predictivo"
% Daniela Alejandra Miranda Ramirez, 2025
%
% Compilar: latexmk -xelatex emi.tex   (desde este directorio)
% =============================================================================

\documentclass{univalle-perfil}

\graphicspath{{../}}
\addbibresource{emi.bib}

% =============================================================================
\begin{document}

% --- Carátula ----------------------------------------------------------------
\universidad{Escuela Militar de Ingeniería}
\facultad{Facultad de Ingeniería de Sistemas}
\carrera{Carrera de Ingeniería de Sistemas}
\tituloTrabajo{Plataforma Web Basada en Inteligencia Artificial y Modelos Predictivos para el Mantenimiento Predictivo de Transformadores Eléctricos: Caso ISI Mustang Bolivia SRL}
\modalidad{Trabajo de Grado}
\tituloOptar{Licenciatura en Ingeniería de Sistemas}
\postulante{Daniela Alejandra Miranda Ramirez}
\tutor{Ing. Marco Antonio Mirabal Condarco}
\ciudad{Santa Cruz}
\pais{Bolivia}
\anio{2025}

\makecaratula

% --- Índice ------------------------------------------------------------------
\pagenumbering{gobble}
\tableofcontents

% --- Contenido ---------------------------------------------------------------
\pagenumbering{arabic}

% =============================================================================
% RESUMEN EJECUTIVO
% =============================================================================
\paginaprevia{RESUMEN EJECUTIVO}

El presente proyecto tiene como finalidad la implementación de una plataforma web
de mantenimiento predictivo para transformadores eléctricos en ISI Mustang Bolivia
SRL, utilizando Inteligencia Artificial (IA) y modelos de Machine Learning. Su
propósito principal es anticipar las necesidades de mantenimiento y optimizar la
gestión operativa mediante la detección temprana de fallos.

La solución se sustenta en una arquitectura de datos tipo Medallion (Bronze, Silver,
Gold), que permite capturar, transformar y gestionar información proveniente de
sistemas como AVEVA PI System, garantizando la calidad, trazabilidad y
disponibilidad de los datos. A partir de esta base sólida, se entrenan modelos
predictivos que analizan variables críticas, como la temperatura del punto caliente,
para identificar patrones de riesgo y generar predicciones confiables.

Se desarrolló un panel interactivo en la web que integra los modelos predictivos con
datos en tiempo real, permitiendo a los usuarios acceder a indicadores clave de
rendimiento, generar alertas tempranas y apoyar la toma de decisiones. Esta
plataforma busca reducir las interrupciones no planificadas, mejorar la eficiencia y
optimizar la gestión del ciclo de vida de los transformadores. Con su implementación,
ISI Mustang Bolivia SRL refuerza su propuesta de valor y se posiciona como líder
regional en mantenimiento predictivo basado en IA para la transformación digital del
sector energético.

\textbf{Palabras clave:} Mantenimiento predictivo, Transformadores eléctricos,
Inteligencia Artificial, Machine Learning, Arquitectura Medallion.

% =============================================================================
% CAPÍTULO I — GENERALIDADES
% =============================================================================
\capitulo{1}{I}{Generalidades}

\subsection{Introducción}

Hoy en día, la transformación digital y la inteligencia artificial son esenciales en la
optimización de procesos industriales y en la toma de decisiones estratégicas. El
mantenimiento de los transformadores eléctricos es esencial para garantizar la
eficiencia operativa y la reducción de costos dentro de la industria energética,
específicamente en la generación de energía eléctrica. Sin embargo, muchas
empresas del sector todavía dependen de enfoques reactivos y no aprovechan el
potencial que ofrecen las técnicas de análisis predictivo, como Machine Learning,
para no solo anticipar un fallo, sino también para mejorar la toma de decisiones
estratégicas.

ISI Mustang, empresa especializada en automatización y digitalización industrial,
utiliza AVEVA PI System, una plataforma para la recolección, almacenamiento y
visualización de datos operativos en tiempo real. Sin embargo, la falta de modelos
predictivos basados en Machine Learning resulta en un comportamiento ineficaz que
impide un mantenimiento predictivo, generando fallas inesperadas y una generación
de energía eléctrica menos eficiente.

\subsection{Antecedentes}

A nivel global, la automatización industrial y la implementación de modelos
predictivos han transformado las operaciones en varias industrias, en particular en el
ámbito energético. Empresas como General Electric (GE), Siemens y Schneider
Electric han sido pioneras en la implementación de soluciones avanzadas de
automatización y análisis predictivo. Según Vijayalakshmi et al. (2023), estas
compañías logran reducir hasta un 40\% los tiempos de inactividad mediante modelos
predictivos basados en IA, lo que impacta directamente en la rentabilidad de sus
operaciones.

En América Latina, países como México y Brasil han sido pioneros en la puesta en
marcha de plataformas de automatización y análisis predictivo. Brasil ha conseguido
disminuir un 15\% las pérdidas técnicas en la distribución de energía mediante la
inteligencia artificial (Keyrus Insights, s.f.).

En Bolivia, aún es nueva la implementación de soluciones avanzadas de inteligencia
artificial en el mantenimiento de infraestructuras energéticas. Las empresas del
sector continúan usando sistemas tradicionales como SCADA y AVEVA PI System,
los cuales permiten la monitorización en tiempo real, pero carecen de capacidades
predictivas avanzadas. ISI Mustang todavía carece de un sistema de análisis
predictivo que posibilite la detección de fallos antes de que sucedan y optimice la
eficiencia del mantenimiento de transformadores eléctricos.

\subsection{Planteamiento del Problema}

\subsubsection{Identificación del problema}

En la actualidad, ISI Mustang y sus clientes del sector energético basan sus
decisiones en los datos proporcionados por AVEVA PI System. Aunque esta
plataforma facilita la recopilación, almacenamiento y visualización de datos
operativos en tiempo real, no cuenta con herramientas de análisis predictivo que
posibiliten la identificación de fallos de forma anticipada.

Las limitaciones de AVEVA PI System restringen el acceso a herramientas avanzadas
de análisis de datos, lo que impide el desarrollo e implementación de estrategias
predictivas. En consecuencia, la empresa depende de enfoques reactivos en lugar de
proactivos para la gestión del mantenimiento. Sin modelos predictivos que analicen
patrones en los datos operacionales y anticipen posibles fallos, se generan
mantenimientos innecesarios o, por el contrario, intervenciones tardías que no
previenen fallos críticos.

\subsubsection{Formulación del problema}

¿De qué forma puede una plataforma web basada en Inteligencia Artificial y modelos
predictivos mejorar la gestión del mantenimiento predictivo de los transformadores
eléctricos en ISI Mustang Bolivia SRL, optimizando la detección temprana de fallos,
la planificación de intervenciones y la eficiencia operativa?

\subsection{Objetivos}

\subsubsection{Objetivo general}

Desarrollar una plataforma web basada en Inteligencia Artificial y modelos predictivos
que optimice el mantenimiento predictivo de los transformadores eléctricos en ISI
Mustang Bolivia SRL, mediante la detección anticipada de fallos, el análisis
inteligente de datos operativos y la automatización de la toma de decisiones
estratégicas para mejorar la eficiencia y reducir costos operativos.

\subsubsection{Objetivos específicos}

\begin{itemize}
    \item Analizar los datos operacionales de ISI Mustang y las necesidades de sus
    clientes en el sector energético.
    \item Diseñar un proceso de Extracción, Transformación y Carga (ETL) para
    integrar datos históricos provenientes de AVEVA PI System que permita
    estructurar y garantizar la calidad de los datos para el análisis predictivo.
    \item Entrenar un modelo predictivo basado en Machine Learning para anticipar
    fallos en los transformadores eléctricos.
    \item Desarrollar una interfaz web interactiva con un dashboard que permita la
    visualización en tiempo real de los datos clave del sector energético, integrando
    un modelo predictivo con Machine Learning.
    \item Evaluar la efectividad de la plataforma mediante pruebas con datos reales y
    métricas de precisión del modelo predictivo, asegurando su correcto
    funcionamiento e integración con los sistemas existentes.
\end{itemize}

\subsection{Justificación}

\subsubsection{Justificación social}

La implementación de esta plataforma consolidará a ISI Mustang como líder en
innovación tecnológica dentro del sector energético boliviano. Al integrar modelos
predictivos y análisis avanzado de datos, la plataforma permitirá optimizar el
mantenimiento predictivo de transformadores eléctricos, reduciendo tiempos de
inactividad, lo que resultará en un suministro eléctrico más estable para la
población. Además, al mejorar la eficiencia de las plantas eléctricas, esto disminuirá
el consumo de combustibles fósiles, reduciendo las emisiones de gases contaminantes.

\subsubsection{Justificación económica}

La implementación de la plataforma de análisis predictivo generará beneficios
económicos significativos para ISI Mustang. Al proporcionar una solución tecnológica
de alto valor agregado, la empresa podrá expandir su cartera de clientes y
diversificar su gama de servicios, lo que resultará en un aumento de sus ingresos.
La automatización de procesos de análisis y predicción reducirá la dependencia de
soluciones externas, permitiendo a ISI Mustang mejorar sus operaciones internas y
reducir costos relacionados con el mantenimiento predictivo de sus clientes del
sector energético.

\subsubsection{Justificación técnica}

La implementación de un sistema de inteligencia artificial con capacidades
predictivas permitirá a ISI Mustang y a sus clientes optimizar la gestión de recursos y
reducir los costos operativos en el sector de energía eléctrica. La implementación de
Machine Learning en el mantenimiento predictivo no solo incrementa la eficiencia
operativa, sino que también reduce los tiempos de inactividad y los costos asociados
a reparaciones no planificadas.

\subsection{Alcances}

\subsubsection{Alcance temático}

El presente trabajo comprende las siguientes áreas de estudio: Metodología de la
Investigación, Estructura de Datos, Programación Avanzada, Base de Datos,
Ingeniería de Software, Ingeniería de Sistemas, Modelación y Simulación de
Sistemas, e Inteligencia Artificial.

\subsubsection{Alcance geográfico}

El presente trabajo se desarrolló en las oficinas de la empresa ISI Mustang SRL,
ubicadas en el 4to Anillo, Centro Empresarial Torre Link, en la ciudad de Santa Cruz
de la Sierra.

\subsubsection{Alcance temporal}

El alcance temporal de este trabajo comprende dos semestres académicos, entre los
meses de febrero y noviembre del año 2025, dando cumplimiento al calendario
académico de la Escuela Militar de Ingeniería.

% =============================================================================
% CAPÍTULO II — MARCO TEÓRICO
% =============================================================================
\capitulo{2}{II}{Marco Teórico}

\subsection{Metodología de la investigación}

El presente proyecto de grado utiliza una investigación de tipo aplicada, ya que busca
generar una solución práctica para un problema concreto dentro del contexto
energético de ISI Mustang. El enfoque es descriptivo-cuantitativo: descriptivo porque
busca especificar las propiedades y características del sistema operativo actual, y
cuantitativo porque se basa en la recolección y análisis de datos operacionales para
entrenar modelos matemáticos de predicción.

\subsection{Inteligencia Artificial y Machine Learning}

La Inteligencia Artificial (IA) es la simulación de procesos de inteligencia humana por
parte de máquinas, especialmente sistemas informáticos. Dentro de la IA, el Machine
Learning (ML) es un subcampo que permite a los sistemas aprender y mejorar
automáticamente a partir de la experiencia sin ser programados explícitamente.

En el contexto del mantenimiento predictivo, los algoritmos de ML analizan patrones
en datos históricos de operación para predecir cuándo un equipo es probable que
falle. Entre los modelos más utilizados para detección de anomalías se encuentran:
el Isolation Forest, los autoencoders LSTM (Long Short-Term Memory) y los modelos
de ensamble. En el presente proyecto se utiliza un modelo ensamble entre un
autoencoder LSTM y un Isolation Forest, ajustado bajo la métrica F2, que logró un
recall del 100\% en la clase NO-NORMAL.

\subsection{Arquitectura Medallion de datos}

La arquitectura Medallion es un patrón de diseño de datos organizado en tres capas:
\textbf{Bronze} (datos crudos en formato inmutable), \textbf{Silver} (datos limpios,
imputados y estandarizados) y \textbf{Gold} (dataset optimizado para consumo
analítico). Este enfoque garantiza trazabilidad, gobernanza de datos y escalabilidad
futura del sistema.

\subsection{Mantenimiento predictivo de transformadores eléctricos}

El mantenimiento predictivo es aquel que se realiza en función del estado real del
equipo, en oposición al mantenimiento preventivo (periódico) y al mantenimiento
correctivo (reactivo). Las variables operativas clave para el monitoreo de
transformadores eléctricos incluyen: temperatura del punto caliente (\textit{hot-spot
temperature}), corriente de carga, tensión de operación y parámetros de aceite
dieléctrico.

\subsection{AVEVA PI System}

AVEVA PI System es una plataforma industrial para la recolección, almacenamiento
y visualización de datos operativos en tiempo real. La PI Web API permite acceder
programáticamente a los datos almacenados en el PI System, lo que facilita la
extracción de series temporales para su posterior procesamiento en el pipeline ETL.

% =============================================================================
% CAPÍTULO III — MARCO PRÁCTICO
% =============================================================================
\capitulo{3}{III}{Marco Práctico}

\subsection{Diseño metodológico}

\subsubsection{Tipo y método de investigación}

En el presente trabajo de grado se utiliza una investigación aplicada orientada a la
investigación de campo. El proceso metodológico siguió las siguientes etapas:
diagnóstico inicial mediante entrevistas al personal técnico y análisis de datos
históricos; recolección de datos operativos desde AVEVA PI System; diseño e
implementación del pipeline ETL con arquitectura Medallion; entrenamiento y
validación del modelo predictivo; y desarrollo e integración de la interfaz web.

\subsubsection{Técnicas e instrumentos de recolección de datos}

Las técnicas utilizadas fueron: entrevistas semiestructuradas con el personal técnico
de ISI Mustang Bolivia SRL, observación directa del funcionamiento del sistema
PI System, y revisión documental de la arquitectura técnica de los transformadores
eléctricos. Estas técnicas permitieron identificar las variables críticas, los
requerimientos funcionales del sistema y las limitaciones del entorno operativo actual.

\subsection{Desarrollo de la solución}

\subsubsection{Pipeline ETL con arquitectura Medallion}

Se implementó un proceso de extracción robusto y modular con gestión segura de
credenciales, que garantizó la captura confiable de datos desde el PI System hacia la
capa Bronze. La capa Silver aplicó limpieza, imputación de valores faltantes y
tratamiento de valores atípicos. La capa Gold generó un dataset optimizado, listo
para el entrenamiento del modelo predictivo.

\subsubsection{Modelo predictivo}

El preprocesamiento incluyó imputación de valores faltantes, normalización y
codificación temporal. Se realizó \textit{feature engineering} construyendo variables
térmicas y eléctricas, además de atributos derivados de ventanas móviles. El modelo
ensamble entre un autoencoder LSTM y un Isolation Forest, ajustado bajo la métrica
F2, logró un recall del 100\% en la clase NO-NORMAL, asegurando que ningún fallo
potencial quedara sin detectar.

\subsubsection{Dashboard web}

Se construyó una interfaz responsiva e interactiva que permite la visualización en
tiempo real de datos operativos y predicciones del modelo. El dashboard, conectado
mediante APIs REST, facilita el monitoreo continuo, desplegando métricas clave y
alertas tempranas a través de indicadores gráficos de fácil interpretación.

\subsection{Pruebas y validación}

Las pruebas con datos históricos y simulaciones en tiempo real confirmaron la
efectividad del pipeline completo. La comparación con el sistema de monitoreo
actual evidenció una mejora sustancial en la capacidad de análisis y toma de
decisiones predictiva, validando que los objetivos del proyecto fueron alcanzados.

% =============================================================================
% CAPÍTULO IV — CONCLUSIONES Y RECOMENDACIONES
% =============================================================================
\capitulo{4}{IV}{Conclusiones y Recomendaciones}

\subsection{Conclusiones}

El desarrollo de la plataforma validó de manera integral la factibilidad de aplicar
arquitecturas modernas de datos y modelos de inteligencia artificial al mantenimiento
predictivo de transformadores eléctricos en ISI Mustang Bolivia SRL.

La adopción de la arquitectura Medallion demostró ser una decisión acertada: la capa
Bronze consolidó datos crudos en formato inmutable, la capa Silver aplicó limpieza,
imputación y tratamiento de valores atípicos, y la capa Gold generó un dataset
optimizado para consumo analítico. Este enfoque aseguró trazabilidad, gobernanza
de datos y bases sólidas para el modelado, además de facilitar la escalabilidad
futura del sistema.

El modelo ensamble entre un autoencoder LSTM y un Isolation Forest, ajustado bajo
la métrica F2, logró un recall del 100\% en la clase NO-NORMAL, asegurando que
ningún fallo potencial quedara sin detectar. Esto validó que tanto el preprocesamiento
como la ingeniería de características fueron determinantes en el rendimiento final del
modelo.

En síntesis, el proyecto alcanzó los objetivos planteados, demostrando que la
combinación de arquitecturas modernas de datos (Medallion), el aprendizaje
automático y la visualización interactiva puede aplicarse de manera efectiva al
mantenimiento predictivo de transformadores. Esto constituye una innovación
tecnológica con impacto directo en la eficiencia operativa, la reducción de costos y
la confiabilidad del sistema eléctrico.

\subsection{Recomendaciones}

Se recomienda mantener un proceso continuo de levantamiento y actualización de
requisitos en coordinación con el cliente, ya que las necesidades pueden evolucionar
conforme se utilice la plataforma en producción. Asimismo, se sugiere extender el
análisis de variables incorporando fuentes externas, como condiciones ambientales,
que podrían enriquecer el modelo predictivo.

Se sugiere evaluar algoritmos adicionales de detección de anomalías, incluyendo
arquitecturas basadas en transformers o ensembles híbridos. En escenarios con
suficientes datos etiquetados, se recomienda explorar modelos supervisados como
XGBoost. Resulta esencial evolucionar hacia un enfoque de MLOps que incorpore el
versionado sistemático de modelos y artefactos, el monitoreo de métricas clave y la
automatización parcial de reentrenamientos para mitigar riesgos de \textit{data drift}
y \textit{concept drift}.

Finalmente, se recomienda consolidar la solución bajo un esquema de MLOps
operativo, que contemple integración continua (CI/CD) para los pipelines de datos y
modelos, monitorización del rendimiento en línea y despliegue controlado de nuevas
versiones del modelo predictivo. Esto asegurará la sostenibilidad y confiabilidad del
sistema a largo plazo.

% =============================================================================
% REFERENCIAS
% =============================================================================
\referencias

\end{document}
